\documentclass{article}

\usepackage[utf8]{inputenc}
\usepackage[T1]{fontenc}
\usepackage[ngerman]{babel}
\usepackage[margin=2.5cm]{geometry}
\usepackage{graphicx}
\usepackage{float}
\usepackage{multirow}
\usepackage{pdflscape}
\usepackage{array}
\usepackage{enumitem}
\usepackage{longtable}

\usepackage{everypage}

\newcommand{\Lpagenumber}{\ifdim\textwidth=\linewidth\else\bgroup
  \dimendef\margin=0 
  \ifodd\value{page}\margin=\oddsidemargin
  \else\margin=\evensidemargin
  \fi
  \raisebox{\dimexpr -\topmargin-\headheight-\headsep-0.5\linewidth}[0pt][0pt]{
    \rlap{\hspace{\dimexpr \margin+\textheight+\footskip}
    \llap{\rotatebox{90}{\thepage}}}}
\egroup\fi}
\AddEverypageHook{\Lpagenumber}

\begin{document}
\pagestyle{empty}

\vspace*{\fill}
\begin{center}
    \Huge
    Anforderungen
\end{center}
\vspace*{\fill}

\newpage

\begin{landscape}

\small
\setlength{\tabcolsep}{6pt}
\renewcommand{\arraystretch}{1.3}

\begin{longtable}{|
    >{\raggedright\arraybackslash}p{2.8cm} |
    >{\raggedright\arraybackslash}p{6.2cm} |
    >{\raggedright\arraybackslash}p{6.2cm} |
    >{\raggedright\arraybackslash}p{5.5cm} |
}
\caption{Anforderungen im Überblick} \\
\hline
 & \textbf{Ursprüngliche Anforderung} & \textbf{Verbesserte Anforderung} & \textbf{User Story} \\
\hline
\endfirsthead

\multicolumn{4}{c}%
{{\bfseries \tablename\ \thetable{} -- Fortsetzung}} \\
\hline
 & \textbf{Ursprüngliche Anforderung} & \textbf{Verbesserte Anforderung} & \textbf{User Story} \\
\hline
\endhead

\hline \multicolumn{4}{r}{{Fortsetzung auf nächster Seite}} \\
\endfoot

\hline
\endlastfoot

\multirow{3}{*}{\textbf{funktional}}
& 1. Nutzerregistrierung und Authentifikation: Erlaubt Nutzern Accounts anzulegen, sich einzuloggen, und ihre Passwörter zurückzusetzen.
& 1.1 Nutzer können Accounts anlegen.
& S1.1 Als Nutzer kann ich einen Account anlegen, sodass ich das System nutzen kann. \\
\cline{3-4}
& 
& 1.2 Nutzer können sich einloggen.
& S1.2 Als Nutzer kann ich mich einloggen, sodass ich auf meine persönlichen Daten und Tasks zugreifen kann. \\
\cline{3-4}
& 
& 1.3 Nutzer können ihre Passwörter zurücksetzen.
& S1.3 Als Nutzer kann ich mein Passwort zurücksetzen, sodass ich wieder Zugriff auf meinen Account erhalte. \\
\cline{2-4}

& 2. User Story- und Backlog-Management: Erlaubt dem Team und speziell dem Product Owner (PO) neue User Stories im Backlog anzulegen. Der PO soll Prioritäten für die User Stories festlegen können.
& 2.1 Nutzer können neue User Stories in einem Backlog anlegen.
& S2.1 Als Nutzer kann ich neue User Stories in einem Backlog anlegen, sodass alle User Stories zentral gesammelt werden können. \\
\cline{3-4}
& 
& 2.2 Der Product Owner kann einer User Story die Priorität \textit{hoch}, \textit{normal} oder \textit{niedrig} zuweisen. 
& S2.2 Als Product Owner kann ich einer User Story die Priorität \textit{hoch}, \textit{normal} oder \textit{niedrig} zuweisen, sodass in der Sprintplanung berücksichtigt werden kann, welche User Stories am dringendsten implementiert werden müssen.\\
\cline{2-4}

& 3. Task Management: Erlaubt dem Team Tasks als Verfeinerung der User Stories zu erstellen (und entsprechend mit den User Stories zu verlinken). Erlaubt dem Team Tasks zu priorisieren und als \textit{complete} oder \textit{incomplete} zu markieren.
& 3.1 Nutzer können Tasks in einem Backlog erstellen.
& S3.1 Als Nutzer kann ich Tasks in einem Backlog erstellen, welche einem Backlog hinzugefügt werden, sodass User Stories in konkrete Arbeitsschritte aufgeteilt werden können. \\
\cline{3-4}
& 
& 3.2 Tasks können mit User Stories verlinkt werden.
& S3.2 Als Nutzer kann ich Tasks mit User Stories verlinken, sodass der Zusammenhang zwischen Tasks und User Stories klar ist. \\
\cline{3-4}
&
& 3.3 Nutzer können einer Task die Priorität \textit{hoch}, \textit{normal} oder \textit{niedrig} zuweisen.
& S.3.3 Als Nutzer kann ich einer Task die Priorität \textit{hoch}, \textit{normal} oder \textit{niedrig} zuweisen, sodass in der Sprintplanung berücksichtigt werden kann, welche Task am dringendsten implementiert werden müssen.\\
\cline{2-4}

& 4. Task/User Story Listen: Erlauben dem Team mehrere Listen für Tasks und User Stories zu managen. Die Listen sollen helfen zwischen Projekt Backlog und Sprintbacklog zu unterscheiden. Zusätzlich soll zwischen folgenden unterschieden werden: Tasks in Bearbeitung, Tasks unter Review, Tasks unter Test, und fertiggestellte Tasks. Das Team soll in der Lage sein, Tasks zwischen den Listen hin- und her- zu bewegen.
& 4.1 Nutzer können Backlogs erstellen.
& S4.1 Als Nutzer kann ich Backlogs erstellen, damit ich die geplanten User Stories und Tasks des jeweiligen Sprints dort sammeln kann.\\
\cline{3-4}
&
& 4.2 Tasks werden in einem Backlog in Tasklisten mit dazugehörigem Status \textit{Wartet}, \textit{In Bearbeitung}, \textit{Review}, \textit{Test} oder \textit{Fertiggestellt} gesammelt.
& S4.2.1 Als Nutzer kann ich in einem Backlog Tasks in Tasklisten mit dazugehörigem Status \textit{Wartet}, \textit{In Bearbeitung}, \textit{Review}, \textit{Test} oder \textit{Fertiggestellt} organisieren, damit alle Nutzer einen Überblick über den aktuellen Bearbeitungsstand einer Task erhalten können.\\
\cline{3-4}
& 
& 4.3 Tasks können von einem Backlog in ein anderes Backlog verschoben werden.
& S4.3 Als Nutzer kann ich Tasks von einem Backlog in ein anderes Backlog verschieben, damit in der Sprint Planung eine Zuordnung von Tasks zu Sprints möglich ist.\\
\cline{3-4}
&
& 4.4 Tasks können von einer Taskliste mit Zustand in eine andere Taskliste mit Zustand verschoben werden.
& S4.4 Als Nutzer kann ich eine Task von einer Taskliste mit Zustand in eine andere Taskliste mit Zustand verschieben, um den Zustand der Task an den aktuellen Bearbeitungsstand anpassen zu können.\\

\cline{2-4}
& 5. Nutzer Profile: Nutzer sollen Profile mit persönlichen Informationen haben (z.B. Name und Rolle (Tester, PO, Entwickler, ...) die Sie ansehen und editieren können.
& 5.1 Nutzer besitzen Profile mit persönlichen Informationen (Name, Rolle, E-Mail, Telefonnummer, Straße, Hausnummer).
& S5.1 Als Nutzer habe ich ein Profil mit persönlichen Informationen (Name, Rolle, E-Mail, Telefonnummer, Straße, Hausnummer), sodass mich andere Teammitglieder identifizieren und kontaktieren können. \\
\cline{3-4}
& 
& 5.3 Nutzer können ihre Profile ansehen.
& S5.3 Als Nutzer kann ich mein Profil ansehen, sodass ich meine gespeicherten Informationen überprüfen kann. \\
\cline{3-4}
& 
& 5.4 Nutzer können den Namen, die E-Mail, die Telefonnummer, die Straße und die Hausnummer ändern.
& S5.4 Als Nutzer kann ich meinen Namen, meine E-Mail, meine Telefonnummer, meine Straße und meine Hausnummer editieren, sodass meine Informationen aktuell bleiben. \\
\cline{2-4}

& 6. Nutzer-Task Zuordnung: Das System erlaubt es eine Zuordnung von Nutzern zu Tasks zu definieren, die auf dem Task-Board zu sehen ist.
& 6.1 Nutzer können Tasks zugeordnet werden.
& S6.1 Als Nutzer kann ich mehrere Nutzer einer Task zuordnen, sodass klar ist, wer woran arbeitet. \\
\cline{3-4}
& 
& 6.2 Die Zuordnung von Nutzern zu Tasks wird im Task-Board angezeigt.
& S6.2 Als Nutzer kann ich die Zuordnung von Nutzern zu Tasks im Task-Board sehen, sodass die Verantwortlichkeiten transparent sind. \\
\cline{2-4}

& 7. Pair-programming Planung: Um Pair-Programming zu planen, soll es möglich sein, mehrere Nutzer zu einer Task zuzuordnen.
& -
& - \\
\cline{2-4}

& 8. Schätzungstracking: Das System erlaubt dem Team Tasks mit Aufwandsschätzungen zu annotieren. Wenn Tasks beendet werden kann zudem annotieret werden wie lange die Lösung der Task tatsächlich gedauert hat. Das Schätzungstracking soll eine Visualisierung bereitstellen die zusammenfasst, wie sehr sich Schätzungen und tatsächliche Bearbeitungszeiten unterscheiden um dem Team die Möglichkeit zu geben, ihre Schätzungen über die Zeit zu verbessern.
& 8.1 Tasks können mit Aufwandsschätzungen versehen werden.
& S8.1 Als Nutzer kann ich Tasks mit Aufwandsschätzungen versehen, sodass der Umfang des Sprints in der Sprint Planung besser geplant werden kann. \\
\cline{3-4}
& 
& 8.2 Die tatsächliche Bearbeitungszeit abgeschlossener Tasks kann erfasst werden.
& S8.2 Als Nutzer kann ich die tatsächliche Bearbeitungszeit erfassen, sodass die Differenz der Schätzung und der tatsächlichen Bearbeitungszeit ermittelt werden kann. \\
\cline{3-4}
& 
& 8.3 Abweichungen zwischen Schätzung und tatsächlicher Bearbeitungszeit werden in einem Balkendiagramm visualisiert.
& S8.3 Als Nutzer kann ich Abweichungen zwischen Schätzung und Realität in einem Balkendiagramm sehen, sodass die Bearbeitungszeit in Zukunft realistischer geschätzt werden kann.\\ 
\cline{2-4}

& 9. Suche und Filter: Das System bietet Such- und Filterfunktionen an um Nutzern zu helfen Tasks zu finden.
& 9.1 Tasks können gesucht werden.
& S9.1 Als Nutzer kann ich Tasks suchen, sodass ich schnell relevante Aufgaben finde. \\
\cline{3-4}
& 
& 9.2 Tasks können gefiltert werden.
& S9.2 Als Nutzer kann ich Tasks filtern, sodass ich nur bestimmte Tasks, welche den Filterkriterien entsprechen, sehe. \\
\cline{2-4}

& 10. Benachrichtigungen: Basierend auf den Schätzungen des Teams, sendet das System eine E-Mail oder in-system Benachrichtigung, um Nutzer an Tasks zu erinnern die bald beendet sein müssen.
& 10.1 Das System sendet interne Benachrichtigungen, welche den Nutzer an Tasks erinnert, die bald beendet werden müssen.
& S10.1 Als Nutzer erhalte ich in-system Benachrichtigungen über Tasks, die bald beendet werden müssen. \\
\cline{2-4}

& 11. Agile Practices: Das Tool soll so viele Agile Practices wie möglich unterstützen.
& -
& - \\
\hline


\multirow{3}{*}{\textbf{nicht funktional}}
& 12. Synchronisation: Neue Angaben die auf einem Rechner von einem Nutzer gemacht wurden sollen in Echtzeit allen anderen Nutzern angezeigt werden.
& 12.1 Änderungen werden allen Nutzern innerhalb von 2 Sekunden angezeigt.
& - \\
\cline{2-4}

& 13. Parallele Bearbeitung: Das Tool soll mehreren Nutzern gleichzeitig erlauben Tasks zwischen den Listen zu bewegen und zu bearbeiten. Konflikte sollen dabei verhindert werden.
& 13.1 Mehrere Nutzer können Tasks gleichzeitig bearbeiten.
& - \\
\cline{3-4}
& 
& 13.2 Tasks können gleichzeitig zwischen Listen bewegt werden.
& - \\
\cline{3-4}
&
& 13.3 Bearbeitungskonflikte werden verhindert.
& - \\
\cline{2-4}

& 14. Multi-Nutzer: Das System soll mehrere Entwickler/Nutzer gleichzeitig handhaben können.
& 14.1 Das System unterstützt 10 gleichzeitige Nutzer.
& -\\
\cline{2-4}

& 15. Verfügbarkeit: Das System soll immer verfügbar sein.
& 15.1 Das System weist eine Verfügbarkeit von 99\% auf.
& -\\
\cline{2-4}

& 16. Updates: Nutzer sollten nicht länger als 2 Sekunden warten um Updates in den Tasklisten zu sehen.
& - 
& - \\
\cline{2-4}

& 17. Änderungen: Das System soll für Entwickler nach dem Release leicht erweiterbar bleiben. Spezifisch sollen Entwickler später einfach neue Tasklistentypen erstellen können, weitere Such und Filtertypen hinzufügen können, weitere Arten von Benachrichtigungen hinzufügen können, und weitere Funktionalität zur Unterstützung von Agile Practices hinzufügen können.
& 17.1 Neue Tasklistentypen können im Quellcode hinzugefügt werden.
& -\\
\cline{3-4}
& 
& 17.2 Neue Suchtypen können im Quellcode hinzugefügt werden.
& -\\
\cline{3-4}
& 
& 17.3 Neue Filtertypen können im Quellcode hinzugefügt werden.
& -\\
\cline{3-4}
& 
& 17.4 Neue Benachrichtigungsarten können im Quellcode hinzugefügt werden.
& -\\
\cline{3-4}
& 
& 17.5 Weitere Funktionen zur Unterstützung agiler Praktiken können im Quellcode ergänzt werden.
& -\\

\end{longtable}
\end{landscape}

\end{document}
